\chapter{The Perverse Engine of the Weil Proof}
\label{chap:perverse-engine}
\providecommand{\pr}{\operatorname{pr}}
\providecommand{\FT}{\operatorname{FT}}

Perverse sheaves are not an ornamental reformulation of the final
cohomology groups.  They perform four operations that ordinary sheaves do
not perform cleanly:
\begin{enumerate}
\item they turn extension across a singular boundary into the canonical
      image \(j_{!*}\) inside an abelian category;
\item they put affine cohomology into one-sided degrees by Artin vanishing;
\item they make the Fourier--Deligne transform an exact equivalence of
      abelian categories;
\item after the weight theorem, they convert upper weight bounds and
      duality into proper purity, intersection-cohomology purity, and the
      decomposition theorem.
\end{enumerate}
This chapter gives the entire theorem chain in the order in which it may be
used.  In particular, purity of \(j_{!*}\) and the decomposition theorem are
\emph{outputs}, not inputs, of the proof of the Weil bounds.

\begin{powerbox}[title={The proof in one display}]
\[
\begin{aligned}
 &j_{!*}+\text{Artin vanishing}
      +\text{Fourier support-forgetting}\\
 &\Longrightarrow
      \FT_\psi:\Perv(\A^1)\xrightarrow{\sim}\Perv(\A^1)\\
 &\Longrightarrow \text{the curve weight estimate}\\
 &\Longrightarrow Rf_!\text{ preserves upper weights}\\
 &\xRightarrow[\ f\text{ proper}\ ]{\text{Verdier duality}}
      Rf_*\text{ preserves purity}\\
 &\Longrightarrow H^i(X_{\Fbar},\Ql)
      \text{ is pure of weight }i\\
 &\Longrightarrow \text{the Riemann hypothesis for }Z(X,T)\\
 &\Longrightarrow
   \begin{gathered}
      j_{!*}\text{ preserves purity},\quad \IC\text{ is pure},\\
      \text{decomposition and relative hard Lefschetz}.
   \end{gathered}
\end{aligned}
\]
\end{powerbox}

\section{The noncircular starting package}

We work with \(E=\Qbarl\), or a finite extension large enough to contain the
eigenvalues under consideration.  The integral category of
\cref{chap:integral-perverse-lattices} supplies lattices and the six
operations; \cref{thm:rationalization-perverse} places us in the self-dual
rational perverse category.  At the beginning of the weight proof the
following facts are available.

\begin{enumerate}
\item The middle perverse \(t\)-structure, its recollement triangles, and
      Verdier self-duality; these are geometric and contain no weight
      theorem.
\item For an open immersion \(j:U\hookrightarrow X\),
      \[
        j_{!*}P=\im\bigl({}^pH^0Rj_!P\longrightarrow{}^pH^0Rj_*P\bigr).
      \]
      It is the unique extension having no nonzero subobject or quotient
      supported on \(X-U\).  These are BBD 1.4.22 and 1.4.25.
\item Every simple perverse sheaf is
      \(j_{!*}(L[\dim U])\) for an irreducible lisse \(L\) on a smooth
      irreducible stratum; this is BBD 4.3.1.
\item Perverse Artin vanishing, BBD 4.1.1--4.1.2.
\item The formal six-operation identities, the trace formula, and local
      monodromy theory, but not the general direct-image weight theorem.
\end{enumerate}

\begin{auditbox}[title={What is not yet available}]
At this stage one may not use
\[
  j_{!*}(\text{pure})\text{ is pure},\qquad
  Rf_!(D_{\leq w})\subset D_{\leq w},\qquad
  Rf_*(\text{pure})\text{ is pure},
\]
or the decomposition theorem.  In BBD these occur in
5.3.2, 5.1.13--5.1.14, and 5.4.5 only after the Deligne weight estimate.
The Fourier proof below uses the \emph{minimality} of \(j_{!*}\), never its
later purity-preservation theorem.
\end{auditbox}

\section{\texorpdfstring{\(j_{!*}\)}{Middle extension}: canonical continuation across a boundary}

\begin{proposition}[The minimal extension mechanism]
\label{prop:jmiddle-engine}
Let \(j:U\hookrightarrow X\) have closed complement \(i:Z\hookrightarrow X\)
and let \(P\in\Perv(U,E)\).  The following characterize \(j_{!*}P\):
\begin{enumerate}
\item \(j^*j_{!*}P\simeq P\);
\item \(j_{!*}P\) has no nonzero perverse subobject and no nonzero perverse
      quotient in the essential image of \(i_*\);
\item \(j_{!*}P\) is the image of the standard-to-costandard morphism
      \({}^pj_!P\to{}^pj_*P\).
\end{enumerate}
It preserves monomorphisms and epimorphisms, though it is not in general an
exact functor on short exact sequences.
\end{proposition}

\begin{proof}
The kernel and cokernel of
\({}^pj_!P\to{}^pj_*P\) are boundary objects by recollement.  Its image
therefore has neither a boundary quotient inherited from the source nor a
boundary subobject inherited from the target.  Conversely, if \(Q\) extends
\(P\) and has the stated minimality, the adjunction maps
\({}^pj_!P\to Q\to{}^pj_*P\) force \(Q\) to be their image.  This proves
uniqueness.  Applying the same boundary argument to the kernel of the image
of a monomorphism and to the cokernel of the image of an epimorphism gives
the final assertion.  See BBD 1.4.22--1.4.26 and
\cref{thm:minimal-extension-characterization}.
\end{proof}

For a smooth dense \(U\) of dimension \(d\), the object
\[
  \IC_X(L):=j_{!*}L[d]
\]
therefore remembers the local system on the top-dimensional stratum and
adds no artificial boundary constituent.  This feature, rather than a
weight assertion, is what is used in the Fourier step.

\section{Perverse Artin vanishing and exact cohomology}

\begin{theorem}[Perverse Artin engine]
\label{thm:perverse-artin-engine}
Let \(a:X\to\Spec k\) be affine and \(P\in\Perv(X,E)\).  Then
\[
  H^r(X_{\bar k},P)=0\quad(r>0),\qquad
  H_c^r(X_{\bar k},P)=0\quad(r<0).
\tag{\thechapter.1}\label{eq:artin-engine-vanishing}
\]
More generally, for an affine \(f:X\to Y\), \(Rf_*\) is right perverse
\(t\)-exact and \(Rf_!\) is left perverse \(t\)-exact.
\end{theorem}

\begin{proof}
The relative assertions are BBD 4.1.1 and its Verdier dual, BBD 4.1.2.
Apply them to \(a\).  Since the perverse structure on a point is the
ordinary one, right \(t\)-exactness of \(Ra_*\) is precisely the first
vanishing, and left \(t\)-exactness of \(Ra_!\) is the second.
\end{proof}

This is stronger than a dimension bound.  If \(U\) is an affine curve and
\(L[1]\) is perverse, then the possible cohomology of \(L\) is centered at
the single perverse degree \(0\).  If invariants and coinvariants eliminate
the two endpoint terms, \(R\Gamma_c(U,L)[1]\) is an honest vector space in
degree \(0\).  Consequently a localization triangle or a Fourier stalk
formula becomes a short exact sequence, not a long complex with
uncontrolled connecting maps.

\begin{proposition}[Perverse Leray filtration]
\label{prop:perverse-leray-engine}
For \(f:X\to Y\) and \(K\in D^b_c(X,E)\) there is a functorial spectral
sequence
\[
 E_2^{a,b}
   =H^a\bigl(Y_{\bar k},{}^pH^b(Rf_*K)\bigr)
 \Longrightarrow H^{a+b}(X_{\bar k},K).
\tag{\thechapter.2}\label{eq:perverse-leray-engine}
\]
It is the spectral sequence of the finite perverse truncation tower.  Before
the decomposition theorem, only the resulting finite Frobenius-stable
filtration may be used.
\end{proposition}

\begin{proof}
Apply \(R\Gamma(Y,-)\) to the exact couple of the perverse truncation tower
of \(Rf_*K\).  Bounded constructibility makes the tower finite.  Every
construction commutes with the Weil structure, so the filtration is
Frobenius-stable.  No degeneration assertion has entered.
\end{proof}

\section{Why Fourier--Deligne is perverse \texorpdfstring{\(t\)}{t}-exact}

Let \(A=\A^1_x\), \(A'=\A^1_{x'}\), let
\(\mu(x,x')=xx'\), and let \(\mathcal L_\psi\) be the Artin--Schreier sheaf
for a nontrivial additive character.  Define
\[
 \FT_\psi(K)
  :=R\pr_{2!}\bigl(\pr_1^*K\otimes\mu^*\mathcal L_\psi\bigr)[1].
\tag{\thechapter.3}\label{eq:fourier-engine-definition}
\]

\begin{theorem}[Fourier--Deligne perverse equivalence]
\label{thm:fourier-perverse-equivalence-engine}
The functor \(\FT_\psi\) is \(t\)-exact for the middle perverse
\(t\)-structure and induces an exact equivalence
\[
  \FT_\psi:\Perv(A,E)\xrightarrow{\sim}\Perv(A',E).
\tag{\thechapter.4}\label{eq:fourier-heart-equivalence}
\]
It sends simple perverse sheaves to simple perverse sheaves.
\end{theorem}

\begin{proof}
Let \(P\) be perverse.  The projection \(\pr_1\) is smooth of relative
dimension \(1\), so \(\pr_1^*P[1]\) is perverse; this is the smooth-pullback
case of BBD 4.2.4--4.2.5.  Tensoring with the lisse rank-one sheaf
\(\mu^*\mathcal L_\psi\) preserves perversity.  The projection \(\pr_2\) is
affine.  Hence \(R\pr_{2!}\) gives the lower perverse bound by BBD 4.1.2,
while \(R\pr_{2*}\) gives the upper bound by BBD 4.1.1.

The decisive input is Laumon's support-forgetting theorem
\cite[Thm.~1.3.1.1]{LaumonFourier}: for this oscillatory kernel the natural
map
\[
 R\pr_{2!}\bigl(\pr_1^*P\otimes\mu^*\mathcal L_\psi\bigr)
 \longrightarrow
 R\pr_{2*}\bigl(\pr_1^*P\otimes\mu^*\mathcal L_\psi\bigr)
\]
is an isomorphism.  The same object therefore satisfies both perverse
bounds and is perverse.  Fourier inversion supplies a quasi-inverse, up to
the standard sign pullback and Tate twist.  Thus the induced functor on the
heart is an exact equivalence and preserves simple objects.  This is exactly
\cite[Thm.~1.3.2.3 and Cor.~1.3.2.4]{LaumonFourier}; see also
\cref{thm:fourier-support-forgetting,thm:fourier-perverse-texact-detailed}.
\end{proof}

\begin{powerbox}
The argument uses perversity three times in an essential way:
\(\pr_1^*[1]\) is exactly normalized, the two affine direct images provide
opposite half-exactness bounds, and support-forgetting identifies those two
images.  Ordinary cohomological degrees do not turn these statements into an
exact equivalence of hearts.
\end{powerbox}

\section{The Laumon curve engine}

We now explain the hard one-dimensional estimate in the form in which it
drives all higher-dimensional weights.  Fix
\(\iota:\Qbarl\hookrightarrow\mathbf C\).

\begin{theorem}[Deligne--Laumon curve theorem]
\label{thm:laumon-curve-engine}
Let \(C_0/\Fq\) be smooth projective, let \(j:U_0\hookrightarrow C_0\) be
dense, and let \(L_0\) be lisse and punctually \(\iota\)-pure of weight
\(w\) on \(U_0\).  Then
\[
  H^r(C_{\Fbar},j_{!*}L[1])
  \quad\text{is pure of weight }w+1+r.
\tag{\thechapter.5}\label{eq:curve-ic-purity-engine}
\]
Equivalently, \(H^1(C_{\Fbar},j_*L)\), or parabolic cohomology in degree
\(1\), is pure of weight \(w+1\).
\end{theorem}

\begin{proof}
We separate the reductions from the Fourier step.

\emph{Step 1: reduction to the affine line.}
After a finite extension of the ground field and a finite map
\(C\to\PP^1\), devissage reduces to an irreducible lisse sheaf \(L\) on a
dense open \(U\subset A=\A^1\), non-geometrically-constant and unramified at
\(\infty\).  The reverse reduction is
\cite[Prop.~4.3.2.1]{LaumonFourier}.  Finite pushforward is perverse
\(t\)-exact by BBD 4.1.3, so the reduction preserves the exact perverse
problem.

\emph{Step 2: make the coefficient into one simple object.}
Put
\[
  P=j_{!*}L[1]\in\Perv(A,E).
\]
Minimal extension and irreducibility make \(P\) simple by BBD 4.3.1.  By
\cref{thm:fourier-perverse-equivalence-engine},
\[
  P'=\FT_\psi(P)
\]
is again simple perverse.  On its maximal lisse open it therefore has the
form \(j'_{!*}L'[1]\) with \(L'\) irreducible.  This is the classification
used in \cite[Thm.~1.4.2.1]{LaumonFourier}.

\emph{Step 3: prove that the Fourier local system is pure.}
The sheaf
\[
  L\oplus L^\vee(-w)
\]
is \(\iota\)-real: every local characteristic polynomial is stable under
complex conjugation.  Grothendieck's trace formula and the Fourier trace
identity show that \(L'\) is an irreducible constituent of an
\(\iota\)-real lisse sheaf on its lisse locus.  Deligne's real-sheaf
criterion \cite[Thm.~1.5.1]{DeligneWeilII} says that every irreducible lisse
constituent of an \(\iota\)-real sheaf on a geometrically connected curve is
pure.  Hence \(L'\) is pure of some weight \(w'\).  Notice that no direct
image weight theorem or purity of \(j_{!*}P\) was used.

\emph{Step 4: read global cohomology as local Fourier monodromy.}
Laumon's stationary-phase theorem identifies the local monodromy of \(L'\)
with a finite filtration whose graded pieces are local Fourier transforms
of the local monodromies of \(L\) at the singular points and at infinity.
After taking the relevant inertia invariants, the global group occurs as a
specified Frobenius-stable subquotient:
\[
  H_c^1(U_{\Fbar},L)
  \quad\text{is a subquotient of a local-monodromy filtration on}\quad
  L'_{\bar\eta_{0'}} .
\tag{\thechapter.6}\label{eq:laumon-stationary-filtration}
\]
In the extra normalization used in Laumon \S4.4---irreducible,
nonconstant input with the prescribed behavior at infinity---this
filtration has a two-step portion whose global graded piece is
\(H_c^1(U_{\Fbar},L)\) and whose neighboring graded piece is the
unramified infinity contribution \(L_{\bar\eta_\infty}(-1)\).  We use this
precise graded-piece statement, not a universal short exact sequence for
arbitrary input.  Perverse Artin vanishing is what concentrates the global
contribution in degree \(1\).  The full filtration and its local summands
are constructed in Laumon 2.3.2.1, 2.4.2.1--2.4.3, and 2.5.3.1; the
specialized global calculation is \cite[\S4.4]{LaumonFourier}.  The later
Fourier and curve chapters spell out the corresponding nearby-cycle and
inertia maps.

\emph{Step 5: determine \(w'\) and bound the global term.}
Deligne's local monodromy estimates
\cite[\S1.8, especially Thm.~1.8.4]{DeligneWeilII} compare the weights of
the inertia-invariant and one-step unipotent graded pieces in
\eqref{eq:laumon-stationary-filtration}.  The normalized infinity piece
\(L_{\bar\eta_\infty}(-1)\) has weight \(w+2\).  Applying the opposite local
inequalities to the two adjacent graded pieces forces \(w'=w+1\).
The global graded subquotient \(H_c^1(U_{\Fbar},L)\) consequently has
weights at most \(w+1\).

\emph{Step 6: turn the upper bound into equality.}
The preceding upper bound is not, by itself, a purity statement for all of
\(H_c^1(U_{\Fbar},L)\): boundary classes may have smaller weights.  Pass
instead to the parabolic image
\[
 H_p^1(U,L):=\im\!\left(H_c^1(U_{\Fbar},L)
             \longrightarrow H^1(U_{\Fbar},L)\right)
 \simeq H^1(C_{\Fbar},j_*L).
\]
Poincare--Verdier duality induces a perfect Frobenius-equivariant pairing
between this image and the corresponding parabolic image for
\(L^\vee(1)\).  Since \(L^\vee(1)\) is pure of weight \(-w-2\), applying
the upper bound to it gives the reciprocal lower bound \(w+1\) on
\(H_p^1(U,L)\).  Thus every Frobenius eigenvalue on parabolic cohomology has
absolute value \(q^{(w+1)/2}\).  The endpoint degrees are immediate from
invariants and coinvariants.  This proves \eqref{eq:curve-ic-purity-engine}.

The result is Deligne \cite[Thm.~3.2.3]{DeligneWeilII} and Laumon
\cite[Thm.~4.1.3]{LaumonFourier}.
\end{proof}

\begin{remark}[Where the genuinely deep curve input sits]
The real-sheaf criterion and the local monodromy theorem are proved before
the general direct-image theorem in Weil II.  Their proofs use rationality
of curve \(L\)-functions, positivity, and local monodromy, but not
Deligne 3.2.3 or 3.3.1.  This placement is the essential noncircularity
check for Laumon's proof.
\end{remark}

\section{From curves to arbitrary \texorpdfstring{\(Rf_!\)}{Rf-shriek}}

For a mixed Weil complex \(K_0\), write \(K_0\in D_{\leq w}\) when
\(\mathcal H^a(K_0)\) has weights at most \(w+a\).  BBD 5.4.1 equivalently
tests this condition on perverse cohomology:
\[
 K_0\in D_{\leq w}
 \Longleftrightarrow
 {}^pH^b(K_0)\text{ has weights at most }w+b.
\tag{\thechapter.7}\label{eq:perverse-weight-test-engine}
\]

\begin{theorem}[Perverse direct-image weight bound]
\label{thm:perverse-direct-image-weight-bound}
Let \(f:X_0\to Y_0\) be separated of finite type.  Then
\[
  K_0\in D_{\leq w}(X_0)
  \Longrightarrow
  Rf_!K_0\in D_{\leq w}(Y_0).
\tag{\thechapter.8}\label{eq:rfshriek-upper-engine}
\]
Equivalently, if \(P_0\) is perverse of weights at most \(w\), then
\[
  {}^pH^b(Rf_!P_0)
  \quad\text{has weights at most }w+b.
\tag{\thechapter.9}\label{eq:perverse-rfshriek-bound-engine}
\]
\end{theorem}

\begin{proof}
The detailed induction is carried out in
\cref{thm:deligne-laumon-direct-image}; here we isolate the role of the
perverse machinery.

First devissage by a smooth stratification reduces to shifted irreducible
lisse sheaves on strata.  Open--closed triangles reduce the complement by
Noetherian induction.  After shrinking the base, relative Noether
normalization factors the remaining map through a quasi-finite morphism to
an affine space over the base.  Quasi-finite fibers are zero-dimensional and
preserve the required bound directly.  Project the affine space one
coordinate at a time.  Each relative \(\A^1\)-projection is controlled
fiberwise by \cref{thm:laumon-curve-engine}.

Perverse Artin vanishing confines the cohomology of every affine projection
to the required half-range.  The perverse Leray spectral sequence
\eqref{eq:perverse-leray-engine} then assembles the coordinate projections
and yields a finite Frobenius-stable filtration.  Weight bounds pass to
extensions and subquotients, so no \(E_2\)-degeneration is needed.  The
closed remainder has smaller dimension and is handled by induction.

Translated back through \eqref{eq:perverse-weight-test-engine}, the result is
\eqref{eq:rfshriek-upper-engine}.  In ordinary cohomological notation this
is Deligne \cite[Thm.~3.3.1]{DeligneWeilII}; the derived/perverse formulation
is BBD 5.1.13--5.1.14 and 5.4.1--5.4.2.
\end{proof}

\begin{corollary}[The dual lower bound]
\label{cor:rfstar-lower-engine}
If \(K_0\in D_{\geq w}(X_0)\), then
\[
  Rf_*K_0\in D_{\geq w}(Y_0).
\tag{\thechapter.10}\label{eq:rfstar-lower-engine}
\]
\end{corollary}

\begin{proof}
By definition \(K_0\in D_{\geq w}\) means
\(\D_XK_0\in D_{\leq-w}\).  Apply
\cref{thm:perverse-direct-image-weight-bound} to \(\D_XK_0\) and use
\[
  \D_YRf_*K_0\simeq Rf_!\D_XK_0.
\]
\end{proof}

\section{Proper purity is a two-line theorem once the engine exists}

\begin{theorem}[Proper purity]
\label{thm:proper-purity-engine}
Let \(f:X_0\to Y_0\) be proper and let \(K_0\) be pure of weight \(w\).
Then \(Rf_*K_0\) is pure of weight \(w\), and
\[
  {}^pH^b(Rf_*K_0)
  \quad\text{is pure of weight }w+b.
\tag{\thechapter.11}\label{eq:proper-perverse-purity-engine}
\]
\end{theorem}

\begin{proof}
The upper bound follows from
\cref{thm:perverse-direct-image-weight-bound} because
\(Rf_*=Rf_!\) for proper \(f\).  The lower bound follows from
\cref{cor:rfstar-lower-engine}.  Their intersection is purity.  Apply the
BBD perverse weight criterion, BBD 5.4.4, to obtain
\eqref{eq:proper-perverse-purity-engine}.
\end{proof}

\begin{powerbox}
This is the central payoff.  The difficult part of the Weil proof is the
one-sided theorem for \(Rf_!\).  Properness identifies \(!\) with \(*\), and
Verdier duality reflects the upper bound into a lower bound.  Exact purity
then appears without a second geometric induction.
\end{powerbox}

\section{The smooth proper Weil theorem}

\begin{theorem}[Cohomological Riemann hypothesis]
\label{thm:weil-rh-perverse-engine}
Let \(X_0/\Fq\) be smooth proper of pure dimension \(d\).  Every eigenvalue
\(\alpha\) of geometric Frobenius on
\(H^i(X_{\Fbar},\Ql)\) satisfies, under every complex embedding,
\[
  |\alpha|=q^{i/2}.
\tag{\thechapter.12}\label{eq:weil-rh-engine}
\]
\end{theorem}

\begin{proof}
The shifted constant sheaf
\[
  P_0=\Ql[d]
\]
is perverse and pure of weight \(d\).  Apply
\cref{thm:proper-purity-engine} to the structure morphism
\(a:X_0\to\Spec\Fq\).  Since perverse and ordinary cohomology agree on a
point, the complex \(Ra_*P_0\) is pure of weight \(d\), so
\[
 H^{i-d}(Ra_*P_0)=H^i(X_{\Fbar},\Ql)
\]
is pure of weight
\[
  d+(i-d)=i.
\]
This is exactly \eqref{eq:weil-rh-engine}.
\end{proof}

Together with the trace formula, the determinant--trace identity, duality,
and the integral lattice argument developed elsewhere in these notes,
\cref{thm:weil-rh-perverse-engine} yields rationality, the functional
equation, integral Frobenius polynomials, independence of \(\ell\), and the
full Weil conjectures.  The perverse calculation explains the otherwise
mysterious degree \(i\): it is the sum of the weight \(d\) of
\(\Ql[d]\) and the residual cohomological degree \(i-d\).

\section{Only now: purity of \texorpdfstring{\(j_{!*}\)}{middle extension} and intersection complexes}

\begin{theorem}[Intermediate extension preserves purity]
\label{thm:jmiddle-purity-engine}
Let \(j:U_0\hookrightarrow X_0\) be a locally closed immersion and let
\(P_0\) be a pure perverse sheaf of weight \(w\) on \(U_0\).  Then
\(j_{!*}P_0\) is pure of weight \(w\).
\end{theorem}

\begin{proof}
For an affine open immersion, \(Rj_!\) and \(Rj_*\) have the complementary
perverse exactness properties of BBD 4.1.1--4.1.3.  The upper direct-image
weight theorem puts the relevant standard extension in weights at most
\(w\); BBD 5.3.1 passes this bound to its image \(j_{!*}P_0\).  Duality and
the minimal-extension identity
\(\D j_{!*}\simeq j_{!*}\D\) provide the lower bound.  A general locally
closed immersion is handled by factoring off a closed immersion and by
Noetherian induction on the boundary.  This is BBD 5.3.2 and 5.3.3.
\end{proof}

\begin{theorem}[Purity of intersection cohomology]
\label{thm:ic-purity-engine}
Let \(X_0\) be irreducible of dimension \(d\), let
\(j:U_0\hookrightarrow X_0\) be a smooth dense open, and let \(L_0\) be
lisse and pure of weight \(u\).  Then
\[
  \IC_{X_0}(L_0)=j_{!*}L_0[d]
\]
is a pure perverse sheaf of weight \(u+d\).  If \(X_0\) is proper, then
\[
  IH^i(X_{\Fbar},L)
    :=H^{i-d}(X_{\Fbar},\IC_X(L))
  \quad\text{is pure of weight }u+i.
\tag{\thechapter.13}\label{eq:intersection-purity-engine}
\]
\end{theorem}

\begin{proof}
The lisse perverse sheaf \(L_0[d]\) is pure of weight \(u+d\).  Apply
\cref{thm:jmiddle-purity-engine}.  If \(X_0\) is proper, apply
\cref{thm:proper-purity-engine} to its structure morphism and compute the
shift:
\((u+d)+(i-d)=u+i\).
\end{proof}

Thus perversity extends the Riemann-hypothesis statement from smooth proper
varieties to intersection cohomology of singular proper varieties.  Ordinary
cohomology of a singular proper variety is generally mixed; \(\IC_X\) is the
canonical pure replacement.

\section{Pure complexes split: the decomposition theorem}

\begin{theorem}[Geometric splitting and semisimplicity]
\label{thm:decomposition-engine}
Let \(K_0\) be a pure complex of weight \(w\) on \(Y_0/\Fq\), and write
\(K\) for its pullback to \(Y=Y_0\otimes\Fbar\).  Then
\[
  K\simeq\bigoplus_b{}^pH^b(K)[-b],
\tag{\thechapter.14}\label{eq:pure-splitting-engine}
\]
and every \({}^pH^b(K)\) is a semisimple perverse sheaf.  Consequently, if
\(f:X_0\to Y_0\) is proper and \(P_0\) is pure perverse, then
\[
 Rf_*P\ \simeq\
 \bigoplus_{b,Z,L}
 \IC_{\bar Z}(L)[-b]
\tag{\thechapter.15}\label{eq:decomposition-engine}
\]
after geometric base change, with semisimple lisse \(L\) on smooth dense
opens of the supports \(Z\).
\end{theorem}

\begin{proof}
By BBD 5.4.4, \({}^pH^b(K_0)\) is pure of weight \(w+b\), so the shifted
object \({}^pH^b(K_0)[-b]\) is pure of weight \(w\).  Let \(a\) be the top
nonzero perverse degree.  The upper truncation triangle has the form
\[
 B_0\longrightarrow K_0\longrightarrow
 A_0:={}^pH^a(K_0)[-a]\xrightarrow{\partial_0}B_0[1],
\]
where induction makes \(B_0={}^p\tau_{\leq a-1}K_0\) pure of weight \(w\).
The connecting morphism is an arithmetic class in
\(\Ext^1(A_0,B_0)\).  BBD 5.1.15(iii) states that its image in geometric
\(\Ext^1\) is zero.  Hence the triangle splits after base change to
\(\Fbar\).  Induction on the finite perverse amplitude proves
\eqref{eq:pure-splitting-engine}.  This is BBD 5.4.5.

For semisimplicity, take a pure perverse \(Q_0\).  Over the geometric base,
choose the maximal semisimple subobject \(S\subset Q\); after a finite
extension it descends.  The quotient \(T\) is pure of the same weight by
BBD 5.3.1.  The extension class of
\(0\to S\to Q\to T\to0\) again maps to zero in geometric Ext by
BBD 5.1.15(iii), so it splits.  Maximality forces \(T=0\).  Thus \(Q\) is
semisimple, which is BBD 5.3.8.  Finally, BBD 4.3.1 identifies its simple
summands as intermediate extensions of irreducible lisse sheaves.

For proper \(f\), \cref{thm:proper-purity-engine} makes \(Rf_*P_0\) pure.
Apply the two conclusions above to obtain
\eqref{eq:decomposition-engine}.
\end{proof}

\begin{warning}[Arithmetic Frobenius need not be semisimple]
\Cref{thm:decomposition-engine} is a geometric splitting after passage to
\(\Fbar\).  It does not say that the Frobenius operator on each arithmetic
cohomology group is semisimple.  Pure Weil sheaves can carry nontrivial
unipotent Frobenius blocks.  BBD 5.3.8 asserts geometric semisimplicity of
the perverse sheaf, not semisimplicity of every chosen Weil structure.
\end{warning}

\begin{corollary}[Perverse Leray degeneracy for pure proper maps]
\label{cor:perverse-leray-degeneration-engine}
Under the proper and pure hypotheses of
\cref{thm:decomposition-engine}, the perverse Leray spectral sequence
\eqref{eq:perverse-leray-engine} degenerates at \(E_2\) after geometric base
change, and its filtration admits a noncanonical splitting.
\end{corollary}

\begin{proof}
Apply \(R\Gamma(Y,-)\) to the direct-sum decomposition
\eqref{eq:pure-splitting-engine}.  The resulting cohomology is the direct sum
of the \(E_2\)-terms, which kills all higher differentials.  Noncanonicity is
inherited from the choices in splitting the perverse truncation triangles.
\end{proof}

If \(f\) is projective and \(\eta\) is the Chern class of a relatively ample
line bundle, BBD 5.4.10 further gives relative hard Lefschetz
\[
  \eta^r:{}^pH^{-r}(Rf_*P_0)
  \xrightarrow{\sim}{}^pH^r(Rf_*P_0)(r)
  \qquad(r\geq0).
\tag{\thechapter.16}\label{eq:relative-hard-lefschetz-engine}
\]
This theorem is another downstream benefit of purity.  It must not be used
to establish the weight theorem from which its proof is derived.

\section{What each perverse tool contributes}

\begin{center}
\begin{tabularx}{\textwidth}{>{\raggedright\arraybackslash}p{0.24\textwidth}
                              >{\raggedright\arraybackslash}X}
\toprule
Perverse construction & Irreplaceable role in the proof\\
\midrule
\(j_{!*}\) &
Packages an irreducible local system as a simple global object without
boundary subobjects or quotients; later produces the pure intersection
complex.\\
\addlinespace
Artin vanishing &
Concentrates affine curve cohomology into the degree in which Fourier
stalks and local monodromy form short exact sequences.\\
\addlinespace
Fourier \(t\)-exactness &
Turns an oscillatory derived integral into an exact equivalence of hearts,
so simplicity survives the transform and Deligne's real-sheaf criterion
applies to one irreducible lisse constituent.\\
\addlinespace
Perverse Leray &
Assembles relative one-dimensional estimates through finite
Frobenius-stable filtrations without assuming premature degeneration.\\
\addlinespace
Perverse weight criterion &
Translates the sheafwise \(Rf_!\) estimate into a statement about exact
perverse constituents with the correct degree shift.\\
\addlinespace
Proper purity &
Combines the \(!\)-upper and \(*\)-lower bounds in one object because
\(!=*\) for a proper map.\\
\addlinespace
Pure splitting &
Kills geometric extension classes between equal-weight truncation pieces,
yielding decomposition, semisimplicity, and perverse Leray degeneration.\\
\bottomrule
\end{tabularx}
\end{center}

\section{Primary-source map and dependency audit}

The exact source map used in this chapter is:
\begin{itemize}
\item BBD 1.4.22, 1.4.25, 1.4.26:
      \(j_{!*}\), minimality, and simple objects under recollement;
\item BBD 4.1.1--4.1.3:
      affine Artin vanishing and quasi-finite \(t\)-exactness;
\item BBD 4.2.4--4.2.5 and 4.2.8:
      smooth shifted pullback and external products;
\item BBD 4.3.1:
      finite length and classification of simple perverse sheaves;
\item Deligne 1.5.1 and 1.8.4:
      the real-sheaf purity criterion and local monodromy weights;
\item Laumon 1.3.1.1, 1.3.2.3--1.3.2.4, and 1.4.2.1:
      support-forgetting, Fourier perverse equivalence, and simple types;
\item Laumon 2.3.2.1, 2.4.2.1--2.4.3, 2.5.3.1, and 4.4:
      local Fourier transform, stationary phase, and the curve proof;
\item Deligne 3.2.3 and Laumon 4.1.3:
      the curve purity theorem;
\item Deligne 3.3.1, BBD 5.1.13--5.1.14 and 5.4.1:
      the general upper direct-image bound and its perverse translation;
\item BBD 5.3.1--5.3.2, 5.3.8, 5.4.4--5.4.5:
      subquotient weights, pure \(j_{!*}\), geometric semisimplicity, and
      pure-complex splitting;
\item BBD 5.4.10:
      relative hard Lefschetz after purity.
\end{itemize}

\begin{auditbox}[title={Acyclic theorem order}]
\[
\begin{array}{c}
\text{pro-etale coefficients and six operations}\\
\downarrow\\
\text{perverse recollement, minimal \(j_{!*}\), Artin vanishing}\\
\downarrow\\
\text{Fourier perverse equivalence and local Fourier transform}\\
\downarrow\\
\text{Deligne--Laumon curve theorem}\\
\downarrow\\
\text{general \(Rf_!\) upper weights}\\
\downarrow\\
\text{proper purity and smooth proper Weil RH}\\
\downarrow\\
\text{pure \(j_{!*}\), IC purity, decomposition, relative Lefschetz}.
\end{array}
\]
Every downward arrow is used; no upward arrow is allowed.  In particular,
decomposition is not a premise of the proof of the Weil conjectures in this
volume.
\end{auditbox}

\section{Formalization interface}

\begin{formalization}
The reusable theorem interfaces suggested by the proof are:
\begin{enumerate}
\item a recollement package whose heart-level image defines \(j_{!*}\) and
      proves the no-boundary universal property;
\item affine left/right perverse exactness as two separate theorems, linked
      by Verdier duality;
\item a Fourier kernel record carrying smoothness of \(\pr_1\), affineness
      of \(\pr_2\), support-forgetting, duality, and inversion;
\item a theorem that those fields imply an exact equivalence of perverse
      hearts, rather than postulating Fourier \(t\)-exactness independently;
\item a curve-engine record containing the real-sheaf criterion, local
      monodromy filtration, the stationary-phase filtration with its specified
      global subquotient, and duality;
\item an upper-weight theorem for \(Rf_!\), with the lower theorem derived
      mechanically by duality;
\item proper purity derived from the equality \(Rf_!=Rf_*\), not stored as
      an unrelated axiom;
\item pure splitting derived from the geometric vanishing of the image of
      the relevant arithmetic \(\Ext^1\)-class.
\end{enumerate}
The final test case should reduce
\(\Ql[d]\) along \(X\to\Spec\Fq\) and normalize the weight by the literal
identity \(d+(i-d)=i\).  This catches simultaneously the most common shift,
twist, and Frobenius-convention errors.
\end{formalization}

